\documentclass[../../../include/open-logic-section]{subfiles}

\begin{document}

\olfileid[pt]{sth}{ordinals}{basic}
\olsection{Propriedades Básicas dos Ordinais}

Observamos que os primeiros ordinais são os números naturais. A razão
principal para desenvolver uma teoria dos ordinais é alargar o
princípio de indução que vale nos números naturais. Iremos construir
isso através duma sucessão de resultados elementares.

\begin{lem}\ollabel{ordmemberord}
Todo o !!{element} de um ordinal é um ordinal.
\end{lem}

\begin{proof}
Seja $\alpha$ um ordinal com $b \in \alpha$. Como $\alpha$ é
transitivo, $b \subseteq \alpha$. Logo $\in$ bem ordena $b$ pois $\in$
bem ordena $\alpha$.

Para ver que $b$ é transitivo, suponha $x \in c \in b$. Logo $c \in \alpha$ pois $b
\subseteq \alpha$. Outra vez, como $\alpha$ é transitivo, $c \subseteq
\alpha$, de modo que $x \in \alpha$. Logo $x, c, b \in \alpha$. Mas $\in$
bem ordena $\alpha$, logo $\in$ é uma relação transitiva em
$\alpha$ por \olref[sth][ordinals][wo]{wo:strictorder}. Logo, como $x
\in c \in b$, temos $x \in b$. Generalizando, $c \subseteq b$
\end{proof}

\begin{cor}\ollabel{ordissetofsmallerord}
$\alpha = \Setabs{\beta \in \alpha}{\beta \text{ é um ordinal}}$, para
qualquer ordinal~$\alpha$
\end{cor}

\begin{proof}
Imediato a partir de \olref{ordmemberord}.
\end{proof}

A ideia grosseira dos dois resultados principais seguintes,
\olref{ordinductionschema} e \olref{ordtrichotomy}, é que os
ordinais eles próprios são bem ordenados por pertença:

\begin{thm}[Indução Transfinita]\ollabel{ordinductionschema}
Para qualquer fórmula $\phi(x)$: 	
\[
	\text{se }\exists \alpha \phi(\alpha)\text{, então }\exists \alpha(\phi(\alpha)
	\land  (\forall \beta \in \alpha) \lnot \phi(\beta))
\]
onde os quantificadores exibidos ficam implicitamente restritos a ordinais.
\end{thm}

\begin{proof}
%\olref{ordinductionschema1}. 
Suponha $\phi(\alpha)$, para algum ordinal $\alpha$. Se $ (\forall \beta
\in \alpha) \lnot \phi(\beta)$, terminamos. Caso contrário, como
$\alpha$ é um ordinal, tem um !!{element} $\in$-mínimo que verifica
$\phi$, e este é um ordinal por \olref{ordmemberord}.
%	\olref{ordinductionschema2}. Suponha que existe algum ordinal
%	$\gamma$ tal que $\lnot\phi(\gamma)$. Então, por
%	\olref{ordinductionschema1}, existe um ordinal $\in$-minimal
%	$\alpha$ para o qual $\lnot\phi(\alpha)$. Logo $(\forall \beta <
%	\alpha) \phi(\beta)$, tornando falso o antecedente do condicional.
\end{proof}
\noindent
Repare que podemos exprimir igualmente \olref{ordinductionschema} como o
esquema:
\[
\text{se }\forall \alpha((\forall \beta \in \alpha)\phi(\beta) \lif 
\phi(\alpha))\text{, então }\forall \alpha\phi(\alpha)
\]
tomando simplesmente $\lnot\phi(\alpha)$ no \olref{ordinductionschema}, e depois fazendo manipulações lógicas elementares.

\begin{thm}[Tricotomia]\ollabel{ordtrichotomy} 
$\alpha \in \beta \lor \alpha = \beta \lor \beta \in \alpha$, para quaisquer
ordinais $\alpha$ e $\beta$. 
\end{thm}

\begin{proof}
A prova é por dupla indução, isto é, usando
\olref{ordinductionschema} duas vezes. Diga-se que $x$ é \emph{comparável}
com $y$ sse $x \in y \lor x = y \lor y \in x$. 

Para a indução, suponha que todo o ordinal em~$\alpha$ é comparável
com \emph{todo} o ordinal. Para uma segunda indução, suponha que
$\alpha$ é comparável com todo o ordinal em~$\beta$. Mostraremos que
$\alpha$ é comparável com~$\beta$. Por indução em~$\beta$, seguirá que
$\alpha$ é comparável com todo o ordinal; e por indução em~$\alpha$,
\emph{todo} o ordinal é comparável com \emph{todo} o ordinal, como
pretendido. Basta supor que $\alpha \notin \beta$ e $\beta \notin
\alpha$, e mostrar que $\alpha = \beta$. 

Para mostrar que $\alpha \subseteq \beta$, fixe $\gamma \in \alpha$; este
é um ordinal por \olref{ordmemberord}. Logo, pela primeira hipótese de
indução, $\gamma$ é comparável com $\beta$. Mas se $\gamma
= \beta$ ou $\beta \in \gamma$ então $\beta \in \alpha$ (invocando a
transitividade de $\alpha$ se necessário), contra a suposição; logo $\gamma \in \beta$. Generalizando, $\alpha \subseteq
\beta$.

Raciocínio exatamente análogo, usando a segunda hipótese de indução,
mostra que $\beta \subseteq \alpha$. Logo $\alpha = \beta$.
\end{proof}\noindent Assim, por vezes escreveremos $\alpha <\beta$
em vez de $\alpha \in \beta$, já que $\in$ se comporta como relação de
ordem. Não há razões profundas para além da familiaridade, e
porque é mais fácil escrever $\alpha \leq \beta$ do que $\alpha \in
\beta \lor \alpha = \beta$.\footnote{Poderíamos escrever $\alpha
\mathrel{\underline{\in}} \beta$; mas isso seria totalmente não
convencional.}

Eis duas consequências rápidas dos últimos resultados, a primeira das
quais põe a nova notação em prática:

\begin{cor}\ollabel{ordordered}
Se $\exists \alpha\phi(\alpha)$, então $\exists \alpha(\phi(\alpha)
\land \forall \beta(\phi(\beta) \lif \alpha \leq \beta))$.
Além disso, para quaisquer ordinais $\alpha, \beta, \gamma$, tanto $\alpha
\notin \alpha$ como $\alpha \in \beta \in \gamma \lif \alpha \in
\gamma$. 
\end{cor}

\begin{proof}
Tal como \olref[wo]{wo:strictorder}.
\end{proof}

\begin{prob}
Complete o raciocínio ``exatamente análogo'' na prova de
\olref[sth][ordinals][basic]{ordtrichotomy}.
\end{prob}

\begin{cor}\ollabel{corordtransitiveord}
$A$ é um ordinal sse $A$ é um conjunto transitivo de ordinais.
\end{cor}

\begin{proof}
\emph{Da esquerda para a direita.} Por \olref{ordmemberord}. \emph{Da direita para a esquerda.}
Se $A$ é um conjunto transitivo de ordinais, então $\in$ bem ordena $A$ por
\olref{ordinductionschema} e \olref{ordtrichotomy}.
\end{proof}

Agora, descrevemos 
\olref{ordinductionschema} e \olref{ordtrichotomy} como dizendo que $\in$ bem ordena os ordinais. Contudo, temos
de ser \emph{muito cautelosos} com este tipo de afirmação, graças ao
resultado seguinte:

\begin{thm}[Paradoxo de Burali-Forti]\ollabel{buraliforti}
Não existe conjunto de todos os ordinais
\end{thm}

\begin{proof}
Por redução ao absurdo, suponha $O$ o conjunto de todos os ordinais. Se $\alpha \in
\beta \in O$, então $\alpha$ é um ordinal, por \olref{ordmemberord}, logo
$\alpha \in O$. Logo $O$ é transitivo, e portanto $O$ é um ordinal por
\olref{corordtransitiveord}. Daí $O \in O$, contradizendo
\olref{ordordered}. 
\end{proof}
\noindent
Este resultado deve o nome a \citeauthor{Burali-Forti1897}. Mas foi Cantor em 1899---numa carta a Dedekind---quem viu
pela primeira vez com clareza a \emph{contradição} em supor que existe um conjunto de
todos os ordinais. Como van
Heijenoort explica:
\begin{quote}
  O próprio Burali-Forti considerou a contradição como estabelecendo,
  por \emph{reductio ad absurdum}, o resultado de que a ordenação natural
  dos ordinais é apenas uma ordenação parcial.
  \citep[p.~105]{Heijenoort1967}
\end{quote}
Deixando de lado o erro de Burali-Forti, podemos resumir o precedente como segue. Os ordinais são conjuntos que individualmente são bem ordenados por
pertença, e coletivamente são bem ordenados por pertença (sem coletivamente constituírem um conjunto). 

Para concluir, eis mais propriedades básicas dos
ordinais que seguem de \olref{ordinductionschema} e
\olref{ordtrichotomy}. 

\begin{prop}
Toda a sucessão estritamente descendente de ordinais é finita.
\end{prop}

\begin{proof}
Toda a sucessão infinita estritamente descendente de ordinais $\alpha_0 > \alpha_1 > \alpha_2 > \ldots$ não tem membro
$<$-mínimo, contradizendo \olref{ordinductionschema}.
\end{proof}

\begin{prop}\ollabel{ordinalsaresubsets}
$\alpha \subseteq \beta \lor \beta \subseteq \alpha$, para quaisquer ordinais
$\alpha, \beta$.
\end{prop}

\begin{proof}
Se $\alpha \in \beta$, então $\alpha \subseteq \beta$ pois $\beta$ é
transitivo. Analogamente, se $\beta \in \alpha$, então $\beta \subseteq
\alpha$. E se $\alpha = \beta$, então $\alpha \subseteq \beta$ e
$\beta \subseteq \alpha$. Logo, por \olref{ordtrichotomy} terminamos.
\end{proof}

\begin{prop}\ollabel{ordisoidentity}
$\alpha = \beta$ sse $\ordeq{\alpha}{\beta}$, para quaisquer ordinais
$\alpha, \beta$.
\end{prop}

\begin{proof}
Os ordinais são boas ordenações; logo isto é imediato da Tricotomia
(\olref[basic]{ordtrichotomy}) e de
\olref[iso]{wellordnotinitial}. 
\end{proof}

\begin{prob}\ollabel{probunionordinalsordinal}
Prove que, se cada membro de $X$ é um ordinal, então $\bigcup X$ é um ordinal.
\end{prob}

\end{document}
